\documentclass[12pt]{article}
\usepackage[utf8]{inputenc}
\usepackage[T1]{fontenc}
\usepackage[margin=1in]{geometry}
\usepackage{setspace}
\usepackage{xcolor}
\usepackage[colorlinks=true, urlcolor=blue!60!black, linkcolor=blue!60!black]{hyperref}
\setlength{\parindent}{0em}
\setlength{\parskip}{0.45em}
\renewcommand{\baselinestretch}{1.0}
\begin{document}
\begin{center}
{\Large \textit{Jones v. Alfred H. Mayer Co.}, 392 U.S. 409 (1968)\par}
\vspace{0.3em}
{\normalsize Case Brief --- One-Page Edition\par}
\end{center}

\textbf{Memory Jogger.} The 1866 Act reaches private housing discrimination under the Thirteenth Amendment (\href{https://supreme.justia.com/cases/federal/us/392/409/}{392 U.S.\ 409}).

\textbf{Facts.} In 1965 a private developer refused to sell Joseph Lee Jones, a Black man, and his white wife a home in its St.\ Louis County subdivision because of his race. They sued under 42 U.S.C.\ \S\ 1982, which guarantees all citizens ``the same right'' to buy property ``as is enjoyed by white citizens.''

\textbf{Procedural History.} The district court dismissed for lack of state action; the Eighth Circuit affirmed.

\textbf{Votes.} Reversed, 7--2 (Stewart for the Court; Douglas concurring; Harlan, joined by White, dissenting).

\textbf{Holding.} Section 1982 bars all racial discrimination, private as well as public, in the sale or rental of property. The Thirteenth Amendment has no state-action limit, and Section 2 lets Congress rationally identify and eliminate the badges and incidents of slavery. \textit{Hodges v.\ United States} was overruled to the extent inconsistent.

\textbf{Reasoning.} The 1866 Congress meant to protect freedmen against private ``custom, or prejudice'' as well as the Black Codes, and the later Fair Housing Act left \S\ 1982 intact. Congress could rationally find that exclusion from white communities, which ``herds men into ghettos,'' was a relic of slavery. Douglas catalogued slavery's persisting badges. Harlan read ``the same right'' as equal legal capacity against state laws and would have dismissed the writ.

\textbf{Examples.} \emph{Same-Side:} a private school excluding Black students (\textit{Runyon}). \emph{Opposite-Side:} a landlord's religious discrimination. \emph{Fence-Sitter:} Section 2 legislation against sexual-orientation discrimination.

\textbf{Critique.} Critics fault the history and warn of tension with \textit{City of Boerne}; defenders say the 1866 Congress targeted private ``white cartels.''

\textbf{Key Quotation.} ``Surely Congress has the power under the Thirteenth Amendment rationally to determine what are the badges and the incidents of slavery, and the authority to translate that determination into effective legislation.'' --- Justice Stewart, 392 U.S.\ at 440.

\end{document}
